\section{Limitations, Future Work and Conclusion}\label{sec:conclusion}

\subsection{Limitations}
\begin{itemize}
  \item \textbf{One shared key per group.} Whoever has it can read every message and could forge a
        sender name; the interface only flags names that do not match the tracker's record.
  \item \textbf{No TLS on HTTP.} The login request from browser to peer and from peer to tracker is
        readable on the network.
  \item \textbf{State in memory.} The tracker forgets who is signed in when it restarts; peers
        re-register automatically but users must sign in again. Message history is kept in memory
        (the last 1000 messages per peer) and is lost when a daemon stops.
  \item \textbf{Simplified CUBIC.} No TCP-friendly region in the strict sense of RFC~8312, no pacing,
        a window of at most 64 chunks, and loss is the only congestion signal.
  \item \textbf{Two fixed channels.} \code{\#general} and \code{\#study}, plus direct conversations;
        no creation of channels, no access control per channel (optional in the brief).
  \item \textbf{One request per HTTP connection} and polling instead of server push; adequate for a
        LAN of a few dozen users.
  \item \todo{anything that did not work in the LAN demo}
\end{itemize}

\subsection{Future work}
\begin{enumerate}
  \item Per-user keys with a key exchange signed by the tracker, so that sender names are proven.
  \item TLS for the HTTP connections (the Python \code{ssl} module wraps the existing sockets).
  \item A delay-based congestion controller (BBR-style pacing) and larger windows for faster links.
  \item Persistent history and user-created channels with membership stored at the tracker.
  \item Persistent TCP connections between peers instead of one connection per message.
\end{enumerate}

\subsection{Conclusion}
The project implements the hybrid model the assignment describes. The tracker does what only a
central server can do --- authenticate users and tell peers where to find each other --- while the
conversation itself and every file go directly from peer to peer, encrypted and authenticated. On
UDP, the application builds its own reliable transport with sequence numbers, selective and
cumulative acknowledgements, a congestion window and fast retransmit, and the measurements show that
it delivers intact files even when a fifth of all datagrams is lost. The design was verified with
\FactTests{} automated tests, most of them against real processes and real sockets.
\todo{one or two sentences on what the team learned}
